\documentclass[11pt]{article}
\usepackage{ochamath}
\subjectcolor{2F5DA8}
\setsheet{線形代数}{対角化　演習}{https://university-math-crj.pages.dev/linear-algebra/diagonalization/}
\namelinetrue
\begin{document}
\makesheethead{問題}

\problem[対角化]
$A=\begin{pmatrix}3&1\\0&2\end{pmatrix}$ を対角化する $P,D$ を1組求めよ。
\workspace{38mm}

\problem[累乗]
問題1の $A$ の $A^2$ を対角化から求め、直接の積と照合せよ。
\workspace{38mm}
\newpage

\problem[対角化できない例]
$J=\begin{pmatrix}2&1\\0&2\end{pmatrix}$ が対角化できない理由を、固有空間から説明せよ。
\workspace{38mm}

\problem[状態遷移]
$M=\begin{pmatrix}0.7&0.2\\0.3&0.8\end{pmatrix}$ とする。$(2,3)^{\mathsf T}$ と $(1,-1)^{\mathsf T}$ の固有値を確認し、初期状態 $(1,0)^{\mathsf T}$ の1回後の確率を求めよ。
\workspace{38mm}

\end{document}
